%% =====================================================================
%%  B4-T-24-01 -- what B4 contributes to "Refusing an Unauthorised Debt"
%%  -------------------------------------------------------------------
%%  LAYER A: APPEND-ONLY. Written once, then left alone. Material found
%%  only here sits in the `unique` slot and is protected by rule R10.
%% =====================================================================
%% @kind: source
%% @id: B4-T-24-01
%% @book: B4
%% @topic: T-24-01
%% @locator: ch. 2, pp. 43--86
%% @status: distilled
%% @unique: yes
%% @integrated: yes
%% @updated: 2026-09-19

\dossier{B4}{T-24-01}{\bookshort{B4}}{ch. 2, pp. 43--86}

\begin{distilled}
  \item The recommended defence against an imposed obligation is to reclassify the initial act rather than to resist the later request: if the benefit was a device, it is not a favour and creates no debt.
  \item Reclassification must happen at the point of the offer, since resisting after acceptance means arguing inside the frame.
  \item Where the benefit cannot be refused, discharging it promptly and in a chosen form removes the balance the technique depends on.
\end{distilled}

\begin{unique}
  \item The reclassification move as a defence --- declining the premise instead of disputing the amount. B1 advises declining unsolicited largeness but does not give the reasoning that makes the refusal holdable under social pressure.
  \item The observation that the defence must be decided in advance because it is hard to execute in front of an audience.
\end{unique}
